This proposition prevents overinterpretation.  A flat eigenvalue trace may be a theorem, not a failed measurement.  Yet the coupling can still affect \(\norm{A^k}\) through interference among open paths.  The synthetic experiment in \cref{sec:synthetic} makes this distinction visible.

\subsection{Gauge transformations and cycle phase}

Let \(D=\diag(e^{i\psi_1},\ldots,e^{i\psi_r})\).  The diagonal unitary similarity
\[
A\mapsto DAD^{-1}
\]
changes an edge phase by \(\psi_i-\psi_j\) but leaves the spectrum unchanged.  Individual edge phases are therefore partly coordinate convention.  The invariant information is accumulated around cycles.

\begin{proposition}[Cycle products are gauge invariant]
For a directed cycle \(i_1\to i_2\to\cdots\to i_m\to i_1\), the product
\[
a_{i_1i_2}a_{i_2i_3}\cdots a_{i_mi_1}
\]
is invariant under diagonal unitary similarity.
\end{proposition}

\begin{proof}
Each transformed factor gains \(e^{i(\psi_{i_k}-\psi_{i_{k+1}})}\).  Multiplying around the cycle telescopes the phases to one.
\end{proof}

Thus a scalar phase sweep is best understood as changing the holonomy of every directed cycle that contains the selected edge.  This is the graph-theoretic reason eigenvalues move.

\section{Why ``spectroscopy'' is the right word}

A spectroscopic experiment applies a controlled variation and reads structure from the response.  The eigenvalue curves are one response, but they require branch tracking and may be ambiguous near collisions.  Trace moments give a branch-invariant alternative.

For \(q\geq1\), define
\[
m_q(\phi)=\tr(A^{(ij)}(\phi)^q).
\]
Since \(m_q\) is the sum of the \(q\)-th powers of the eigenvalues, it is spectral but does not require individual labels.

\begin{theorem}[Trace-moment harmonics count closed walks]
Let
\[
m_q(\phi)=\sum_{\ell=0}^{q}c_{q,\ell}e^{i\ell\phi}.
\]
Then \(c_{q,\ell}\) is the total baseline weight of length-\(q\) closed walks that traverse the selected directed edge \(i\to j\) exactly \(\ell\) times, with the baseline phase of \(a_{ij}\) included in that weight.
\end{theorem}

\begin{proof}
The trace of \(A^q\) is the sum over closed index sequences
\[
i_0\to i_1\to\cdots\to i_q=i_0
\]
of products \(a_{i_0i_1}\cdots a_{i_{q-1}i_q}\).  In the phase family, every occurrence of the selected edge contributes a factor \(e^{i\phi}\).  A closed walk using that edge \(\ell\) times therefore contributes to the \(\ell\)-th Fourier harmonic.  Summing like powers of \(e^{i\phi}\) gives the result.
\end{proof}

This theorem supplies a concrete interpretation of the signal.  Low-order harmonics expose short feedback loops; higher harmonics expose repeated circulation.  In optimizer language, these are recurrent routes through parameter, momentum, and adaptive-state couplings.

\begin{remark}
The phase-moment spectrum remains basis-dependent under arbitrary similarity because the selected entry changes meaning.  Once the perturbation family is fixed, however, the trace response itself is invariant to eigenvalue labeling and numerically stable near branch exchange.
\end{remark}

\section{Projection without materializing the Jacobian}

The full optimizer-state Jacobian is too large for foundation-scale systems.  A diagnostic must therefore operate through products \(v\mapsto J_tv\).

\subsection{Functional optimizer maps}

The cleanest implementation makes the optimizer step a pure function
\[
F_t(z;\mathcal B_t,\xi_t,h),
\]
where \(\mathcal B_t\) is the batch, \(\xi_t\) is explicit randomness, and \(h\) contains hyperparameters.  Freezing \((\mathcal B_t,\xi_t,h)\) makes a Jacobian-vector product well-defined.  Automatic differentiation can evaluate it directly; centered finite differences provide a verification path.

Hessian-vector products have long made curvature information available without materializing the Hessian \cite{pearlmutter1994}.  The same principle applies to the complete optimizer map.

\subsection{Petrov--Galerkin reduction}

Choose right and left bases \(V,W\in\C^{n\times r}\) with \(W^*V=I_r\).  The reduced operator is
\begin{equation}
\widehat J=W^*JV.
\label{eq:petrov}
\end{equation}
A one-sided orthonormal basis uses \(W=V=Q\).  For strongly non-normal operators, two-sided projection is preferable because right and left dynamical subspaces may differ substantially.

Useful basis sources include:
\begin{itemize}[leftmargin=1.5em]
\item Krylov vectors \(v,Jv,J^2v,\ldots\), generated by Arnoldi;
\item recent gradients and optimizer updates;
\item block summaries aligned with layers or parameter tensors;
\item approximate right and left singular vectors of short-horizon propagators;
\item random sketches for inexpensive coverage.
\end{itemize}
Krylov methods have already been used to construct practical low-dimensional optimization subspaces in deep learning \cite{vinyals2011}.  Here they serve measurement rather than direct descent.

Every reduced spectral claim must include Ritz residuals, basis orthogonality error, rank sensitivity, and probe seed.  Nonnormal projection can look persuasive while being wrong; the residual is the receipt.

\subsection{Frozen-step and short-horizon operators}

The frozen Jacobian is the first object, not the final one.  Training is nonautonomous:
\[
